% !TEX root = ../project_report.tex
\section{Predicting the next 30 minutes}
\subsection{Prediction target and available information}
For station $s$, issue time $t$ and threshold $q$ in millimetres, the target is
\[
Y_{s,t}=\sum_{j=1}^{6}r_{s,t+5j\,\mathrm{min}},\qquad
p_{s,t}(q)=P(Y_{s,t}>q\mid X_{s,t}).
\]
The output is the probability of exceeding the requested amount. Observation timestamps mark the assumed endpoints of five-minute accumulation intervals. Labels require all six exact future observations; totals are rounded to 0.001 mm and exceedance is strictly greater than $q$. Examples with missing future observations are excluded.

\textbf{A (short-term forecasting)} uses recent station rainfall as well as station and calendar features. \textbf{B (calendar risk)} uses only station and calendar features to estimate historical risk at a requested time. We fit the same model families in both settings to compare short-term forecasting with calendar risk. The query's three user inputs are station, issue time and threshold; A derives additional features internally from the supplied observation archive.

A reads a supplied clean observation table and requires twelve history endpoints at $t-65,t-60,\ldots,t-10$ minutes. Both recorded update timestamps for every history record must be no later than $t$. The ten-minute buffer reflects the typical publication delay (Table~\ref{tab:updatedelay}); the record-level checks also exclude observations with longer delays. Missing or unavailable history rejects the example or query. Figure~\ref{fig:t2timeline} separates this history from the future target.

\begin{table}[htbp]\centering\small
\begin{tabular}{lrrrr}\toprule
Period & Observations & Median (min) & P90 (min) & Within 10 min (\%)\\\midrule
2017--2021 (training) & 28,256,245 & 9.58 & 15.08 & 68.93 \\
2017--2024 (all years) & 48,538,625 & 9.47 & 12.68 & 76.40 \\
\bottomrule\end{tabular}
\caption{Observation-to-update delay in the cleaned records, using the later of the two update timestamps. P90 is the 90th percentile; the final column gives the percentage updated within ten minutes.}
\label{tab:updatedelay}
\end{table}

\plot{17_prediction_timeline}{The information used by A at issue time $t$. Markers show interval endpoints. The latest history interval ends ten minutes before issue time; the six intervals after issue time form the target.}{fig:t2timeline}
\FloatBarrier
\subsection{Features and their purpose}
Location and station ID represent spatial differences; cyclic time encodings represent the seasonal and daily patterns explored in Part 1 without a discontinuity at midnight or New Year. Annual phase accounts for leap years, and numeric year is excluded. A's rainfall summaries represent recent intensity, persistence and change; all history windows end at $t-10$ minutes (Table~\ref{tab:input_features}).
\begin{table}[htbp]\centering\small
\begin{tabular}{p{.63\linewidth}cc}\toprule
Feature group (number of numeric features) & A & B\\\midrule
Latitude and longitude (2) & Yes & Yes\\
Time of day: sine and cosine (2) & Yes & Yes\\
Annual phase: sine and cosine (2) & Yes & Yes\\
Latest three five-minute rainfall amounts (3) & Yes & No\\
Rainfall totals over 15, 30 and 60 minutes (3) & Yes & No\\
Maximum five-minute rainfall within 30 minutes (1) & Yes & No\\
Wet-interval count and dry-interval fraction over 60 minutes (2) & Yes & No\\
Difference between the latest and preceding 15-minute totals (1) & Yes & No\\\midrule
\textbf{Total numeric features} & \textbf{16} & \textbf{6}\\
Station ID (categorical; counted separately) & Yes & Yes\\\bottomrule
\end{tabular}
\caption{Input features for the two information settings. A uses recent rainfall in addition to the station and calendar inputs shared with B. For LR and GBT, the query threshold $q$ is an additional input; RF uses it when converting class probabilities to exceedance probabilities.}
\label{tab:input_features}
\end{table}

\FloatBarrier
\subsection{Models: answering a query at any rainfall threshold}
We compare a historical-frequency baseline with three Spark MLlib classifiers~\cite{sparkml}. The key modelling choice is how to turn their outputs into $P(Y>q)$ for a user-supplied threshold, including values absent from the fitting grid.

\textbf{Historical frequency} estimates the exceedance rate for each station/month/hour group, shrinking it towards the global training rate using 100 pseudo-observations. Unseen groups use the global rate. This establishes how much can be achieved from recurring rainfall patterns before adding a classifier.

\textbf{Distributional random forest (RF)} predicts seven amount bins: $\{0\}$, $(0,1]$, $(1,2]$, $(2,5]$, $(5,10]$, $(10,20]$ and $(20,\infty)$. Predicting a distribution lets one fitted forest answer different rainfall thresholds. For a query inside a bin, we use the observed training distribution within that bin:

\[
\widehat p(q\mid X)=\sum_{k=0}^{6}\widetilde\pi_k(X)\widehat S_k(q),
\qquad
\widetilde\pi_k=\frac{\pi_k^{1/T}}{\sum_\ell\pi_\ell^{1/T}}.
\]
Here $\pi_k(X)$ is the forest probability for bin $k$, and $\widehat S_k(q)$ is the fraction of training amounts in that bin exceeding $q$. The temperature $T$ adjusts how concentrated the bin probabilities are; $T=1$ leaves them unchanged. The resulting exceedance probability is non-increasing in the requested rainfall threshold. The within-bin distributions are shared across weather conditions and extend to the largest training total, 71.8 mm; the model returns zero at or above this empirical maximum. This is a finite-support limitation, not a claim that larger future accumulations are impossible. The query output flags thresholds outside the training support.

Station IDs are encoded as categorical features, and numeric imputation uses training-set medians. The forests in the main model/fraction comparison use 40 trees, depth 8, a minimum of 25 examples per child node, square-root feature selection, subsampling 0.8 and at most 90 feature bins. The 15\% and 100\% searches select depth 8 from 5 and 8; the 30\% comparison reuses the same reference configuration. Temperature selection chooses $T=1$ for all three fractions in both settings on the 2023 calibration set.

\textbf{Threshold-conditioned logistic regression (LR)} instead treats $q$ as an input and predicts the binary label $\mathbf1(Y>q)$. Each sampled station-time example is expanded to seven thresholds: 0, 0.2, 1, 2, 5, 10 and 20 mm. Station IDs are one-hot encoded and numeric inputs scaled. L2 regularisation is selected from 0.01 and 0.1, with at most 100 optimisation iterations. LR supplies a simple direct-probability comparator and receives no separate calibration fit.

\textbf{Threshold-conditioned gradient-boosted trees (GBT).} GBT uses the same binary labels and threshold expansion as LR, but learns nonlinear relationships by fitting trees sequentially. Its pipeline indexes station IDs, imputes numeric features and assembles them with $q$. We compare depths 3 and 5 with 40 boosting iterations, step size 0.1, a minimum of 25 training examples per child node, subsampling 0.8 and a maximum of 128 feature bins. Every reported configuration selects depth 5. Temperature calibration is applied to the binary probability pair $(1-p,p)$: A selects $T=0.75$ and B selects $T=1$ at each fraction. Unlike the RF survival mixture, this construction does not enforce decreasing probability as $q$ increases.

\subsection{Evaluation protocol}
Adjacent examples share history and future observations, so a random train/test split would mix closely related rainfall episodes. We use chronological partitions (Table~\ref{tab:t2splits}) and remove labels whose target end reaches the following split boundary.
\begin{table}[htbp]\centering\small
\begin{tabular}{lrrp{.29\linewidth}}\toprule
Period & A examples & B examples & Purpose\\\midrule
2017--2021 & 18,688,263 & 27,434,001 & Training population \\
2022 & 6,133,684 & 7,107,898 & Model and cutoff selection \\
2023 & 5,465,591 & 6,719,274 & RF and GBT temperature selection \\
2024 & 5,729,508 & 6,374,476 & Model evaluation \\
\bottomrule\end{tabular}
\caption{Eligible station-time examples after completeness checks and boundary purging. A additionally requires complete published history.}\label{tab:t2splits}
\end{table}

For each classifier in each setting, we compare 15\%, 30\% and 100\% of the eligible training anchors. A deterministic station--time hash (seed 2026) gives nested, label-independent samples. The 30\% samples contain 5,608,769 anchors for A and 8,233,882 for B. Historical frequencies and RF within-bin distributions always use all eligible training anchors; only classifier fitting is subsampled. Selection, calibration and evaluation use the full eligible populations within each setting. Thus the experiment measures classifier fitting-fraction sensitivity; a 15\% fit still uses the full training population for the empirical supports and historical frequencies.

Our primary criterion is mean Brier across the seven fitting thresholds, giving each threshold equal weight. At a particular threshold,
\[
\mathrm{BS}=\frac{1}{N}\sum_{i=1}^{N}(\widehat p_i-z_i)^2,
\qquad z_i=\mathbf1(Y_i>q),
\]
which measures squared probability error; lower is better. We prioritise this score because the task requests probabilities. PR-AUC and ROC-AUC measure event ranking, log loss penalises confident errors, and reliability compares predicted probabilities with observed frequencies. Spark computes AUC with 1,000-bin curve downsampling~\cite{sparkeval}. We examine baseline skill and high-threshold calibration alongside mean Brier to assess performance on rare heavy-rainfall events.

Hyperparameters are selected by 2022 mean Brier. RF/GBT temperature is selected on 2023 from 0.5, 0.75, 1, 1.25, 1.5 and 2; alarm cutoffs use 2022 after applying the selected temperature. RF30 reuses saved reference-configuration weights before calibration and testing. The model/fraction comparison uses 2024 throughout. Later full-training and GBT experiments were developed after the initial 2024 results were inspected. Although 2024 is excluded from fitting and temperature selection, these later comparisons and the final model choice are exploratory rather than an untouched confirmatory test. Execution history and available artifacts are documented in Appendix~\ref{app:experiments}.

\clearpage
\section{Prediction results}
We first assess the value of recent rainfall, then compare model families and training fractions. Calibration and threshold diagnostics explain the remaining prediction errors.

\subsection{Recent rainfall improves forecasts, with smaller gains for heavy rain}

Comparing the 15\% RFs on the same 5,729,508 test examples, A reduces mean Brier from 0.02283 to 0.01613, a \textbf{29.4\% relative reduction} against B. The shared cases cover 89.9\% of B's eligible 2024 set; each forest is fitted on its own eligible training population.
\par
\begin{table}[!ht]\centering\small
\begin{tabular}{rrrrrr}\toprule
$q$ (mm) & Event rate (\%) & Brier A & Brier B & PR-AUC A & PR-AUC B\\\midrule
0 & 6.398 & 0.03798 & 0.05924 & 0.6091 & 0.1097 \\
0.2 & 4.892 & 0.03090 & 0.04614 & 0.5582 & 0.0881 \\
1 & 2.579 & 0.01870 & 0.02499 & 0.4215 & 0.0484 \\
2 & 1.686 & 0.01349 & 0.01652 & 0.3155 & 0.0329 \\
5 & 0.829 & 0.00732 & 0.00820 & 0.1885 & 0.0176 \\
10 & 0.378 & 0.00354 & 0.00376 & 0.1020 & 0.0085 \\
20 & 0.100 & 0.00098 & 0.00100 & 0.0304 & 0.0029 \\
\bottomrule
\end{tabular}
\caption{The 15\% RFs evaluated on identical 2024 examples, with event frequency at each rainfall threshold.}\label{tab:t2common}
\end{table}
\par
The benefit is greatest for any positive rain (35.9\% Brier reduction) and falls to 1.7\% above 20 mm (Table~\ref{tab:t2common}, Figure~\ref{fig:t2common}). Recent history therefore adds useful predictive information in this holdout, but its strongest benefit is at lower rainfall thresholds.
\plot{18_common_comparison}{The 15\% RF comparison on shared 2024 examples. Left: relative Brier reduction with conditional 95\% block-bootstrap intervals. Right: PR-AUC and event prevalence on a logarithmic scale. Rainfall thresholds are displayed as equally spaced categories.}{fig:t2common}
To account for overlapping labels and shared weather across stations, we form 2,000 bootstrap replicates by resampling circular blocks of seven consecutive calendar days, retaining all stations and thresholds together. The 95\% interval for the mean relative Brier reduction is 25.4--32.9\%, conditional on these fitted models and the observed 2024 network.

Equal threshold weights do not equalise the attainable absolute error reduction. The 0 and 0.2~mm queries account for 77.8\% of A's total absolute Brier improvement over B. At 20~mm there are 5,715 positive, overlapping station-time windows: A's Brier is 0.0009791 versus 0.0009975 for always predicting zero probability, only 1.84\% lower; its skill over historical frequency is 2.18\%. Its PR-AUC of 0.0304 still exceeds the event prevalence of about 0.001. Heavy-rain ranking is useful, while the probability-error gain is modest.

\FloatBarrier
\subsection{Model comparison and uncertainty}
Holding the fitting fraction at 100\%, \textbf{RF has the lowest observed mean Brier among the tested classifier pipelines in both settings}. In B, however, the historical-frequency baseline has lower error than all three classifiers (Table~\ref{tab:t2baseline}).
\par
\begin{table}[!ht]\centering\small
\begin{tabular}{llrrrr}\toprule
Mode & Model & Training & Mean Brier & \shortstack{PR-AUC\\at 1 mm} & \shortstack{Log loss\\at 1 mm}\\\midrule
A & Random forest & 100\% & 0.01614 & 0.4229 & 0.0824 \\
A & Logistic regression & 100\% & 0.01944 & 0.3634 & 0.0875 \\
A & GBT & 100\% & 0.01623 & 0.4282 & 0.0824 \\
A & Historical frequency & 100\% & 0.02270 & 0.0467 & 0.1166 \\
B & Random forest & 100\% & 0.02275 & 0.0518 & 0.1167 \\
B & Logistic regression & 100\% & 0.02301 & 0.0428 & 0.1165 \\
B & GBT & 100\% & 0.02288 & 0.0522 & 0.1194 \\
B & Historical frequency & 100\% & 0.02252 & 0.0516 & 0.1146 \\
\bottomrule
\end{tabular}
\caption{Full-training models evaluated on each mode's eligible 2024 population: 5,729,508 examples for A and 6,374,476 for B. Mean Brier averages the seven original thresholds.}\label{tab:t2baseline}
\end{table}
\par
The ranking metrics explain why model choice depends on the objective. In A, GBT has slightly higher PR-AUC at 1 mm (0.4282 versus 0.4229), but higher mean probability squared error. LR has higher ROC-AUC at 1 mm (0.8611 versus 0.8036), while RF has better PR-AUC, Brier and log loss. Ranking events and assigning accurate probabilities therefore describe different aspects of performance.

The paired uncertainty analysis in Table~\ref{tab:pairedmodels} supports a small RF advantage over GBT in A. In B, the mean difference is less clear: the interval includes zero, and GBT has lower Brier at the two lowest thresholds, whereas RF performs better at higher thresholds. Historical frequency has the lowest mean Brier in B; the full RF's mean error is 1.04\% higher, with conditional interval 0.22--1.79\%. This favours the simpler calendar baseline at the evaluated thresholds.

\begin{table}[htbp]\centering\small
\begin{tabular}{llrr}\toprule
Candidate versus reference & Mode & Reduction (\%) & Conditional 95\% interval\\\midrule
RF 100\% versus GBT 100\% & A & +0.560 & [0.222, 0.923] \\
RF 100\% versus GBT 100\% & B & +0.571 & [-0.229, 1.605] \\
RF 100\% versus RF 15\% & A & -0.036 & [-0.092, 0.019] \\
RF 100\% versus RF 15\% & B & -0.047 & [-0.133, 0.021] \\
\bottomrule\end{tabular}
\caption{Relative reduction in mean 2024 Brier, $100(1-\mathrm{BS}_{candidate}/\mathrm{BS}_{reference})$. Point estimates use observed 2024 scores; 95\% intervals use 2,000 paired circular seven-day block resamples. Positive values favour the candidate. Fractions refer to classifier fitting anchors.}\label{tab:pairedmodels}
\end{table}

\FloatBarrier
\subsection{Larger fitting samples bring little improvement at these settings}
The three-fraction comparison examines sensitivity to the number of classifier fitting anchors at the reported settings. \textbf{The recorded mean Brier changes from 15\% to 100\% are below 0.1\% in magnitude for every classifier and setting.}
\begin{table}[!ht]\centering\small
\begin{tabular}{llrrrr}\toprule
Mode & Model & 15\% & 30\% & 100\% & Improvement (\%)\\\midrule
A & Random forest & 0.01613051 & 0.01613151 & 0.01613638 & -0.036 \\
A & Logistic regression & 0.01944150 & 0.01943988 & 0.01944300 & -0.008 \\
A & GBT & 0.01622929 & 0.01623889 & 0.01622724 & +0.013 \\
B & Random forest & 0.02274285 & 0.02274079 & 0.02275362 & -0.047 \\
B & Logistic regression & 0.02301006 & 0.02301067 & 0.02301076 & -0.003 \\
B & GBT & 0.02290321 & 0.02290187 & 0.02288430 & +0.083 \\
\bottomrule\end{tabular}
\caption{Mean 2024 Brier over seven thresholds on each mode's eligible population. Improvement is the relative Brier reduction from 15\% to 100\%; positive means lower error. The 15\% A GBT value comes from a retained aggregate summary; other cells use metric exports.}\label{tab:trainingsampling}
\end{table}
For RF, the 100\% score is 0.036\% higher in A and 0.047\% higher in B than at 15\%. LR also changes little, while GBT improves slightly at full training. A's 30\% GBT score is worse than its 15\% score, so the observed learning curve is not monotonic. The model ordering is unchanged: RF has lower mean Brier than LR and GBT at every fitting fraction in both settings. The paired RF intervals in Table~\ref{tab:pairedmodels} include zero for both settings. These observations apply to one sampling seed, the tested model capacities and the threshold-average objective; they do not establish equivalence at every rainfall threshold or an end-to-end computing-cost reduction. The detailed learning-curve figure is retained in Appendix~\ref{app:experiments}.

The controlled RF comparison on 2022 provides a complementary check. With fitting settings held fixed, increasing the sample from 15\% to 30\% changes mean Brier by only 0.012\% in A's favour and 0.003\% against B; both conditional intervals include zero. A modest increase in tree depth is more useful for A than doubling the sample at the reference settings. Appendix~\ref{sec:rfsensitivity} gives the individual settings and paired intervals.

Even the 15\% samples contain millions of anchors, many with overlapping windows from the same rainfall episode, but only 2,886 (A) and 4,198 (B) totals above 20 mm.

\FloatBarrier
\subsection{Reference-forest calibration and the alarm tradeoff}
For the 15\% RFs at $q=1$ mm, A distinguishes a wider range of risks than B. Mean predicted risk is 2.38\%, against 2.58\% observed, and the ten-bin weighted absolute calibration error is 0.196 percentage points. Most examples occupy low-probability bins; several higher-probability bins underestimate the observed event frequency (Figure~\ref{fig:t2reliability}).
\plot{20_reliability}{Calibration of the 15\% RFs at 1 mm on shared examples. Upper panels compare bin-average predictions with observed event rates; B is enlarged to show its two occupied bins. Lower panels show sample counts in fixed 0.1-wide probability bins across the full probability range.}{fig:t2reliability}

The rainfall threshold defines the event, while the probability cutoff determines whether to issue an alarm. For the 15\% A RF at $q=1$ mm, changing the probability cutoff from 0.5 to the validation-selected 0.2 raises recall from 24.7\% to 46.4\%, while precision falls from 69.3\% to 47.4\%. Alarm cutoffs maximise F1 on 2022 data after applying the temperature selected on 2023. An always-negative rule already achieves about 97.4\% accuracy here, which is why probability quality and the precision--recall tradeoff are more informative than accuracy.

\FloatBarrier
\subsection{Threshold diagnostics: mean-risk bias and unseen queries}
This comparison uses RF, LR and GBT all fitted at 30\%, evaluated on the same eligible population within each setting. It tests two requirements of the probability interface: plausible risk at large rainfall thresholds, and useful predictions between the fitted query values.

\textbf{Large-accumulation risk.} At the matched 30\% fitting fraction, GBT's mean predicted risk above 20 mm is 0.830\% in A and 2.645\% in B, compared with observed frequencies of 0.100\% and 0.103\%. This overprediction helps explain why event ranking and probability error can favour different models. Figure~\ref{fig:t2thresholdrisk} assesses mean-risk bias: agreement of the overall mean with prevalence alone does not establish calibration across individual risk levels. The retained 15\% RF reliability tables put 99.983\% of A's 20~mm cases and all of B's cases in the first 0.1-wide bin. Finer low-risk reliability bins with event counts and block intervals require rescoring individual predictions; they cannot be reconstructed from these coarse aggregates.
\plot{23_threshold_risk}{Observed frequencies and mean predicted probabilities for matched 30\% models at the seven fitting thresholds, on each mode's eligible 2024 examples. The vertical axis is logarithmic and thresholds are equally spaced categories.}{fig:t2thresholdrisk}

\textbf{Queries between fitted thresholds.} The additional queries at 3, 7.5 and 15 mm fall inside RF amount bins and between the values used to fit LR and GBT. RF has lower Brier at all three query values in both settings (Table~\ref{tab:interior}), supporting its within-bin conversion at these evaluated points.
\begin{table}[!ht]\centering\small
\begin{tabular}{lrrrr}\toprule
Mode & $q$ (mm) & RF Brier & LR Brier & GBT Brier\\\midrule
A & 3 & 0.010546 & 0.012359 & 0.010717 \\
A & 7.5 & 0.005051 & 0.006950 & 0.005132 \\
A & 15 & 0.001854 & 0.003118 & 0.001913 \\
B & 3 & 0.012385 & 0.012510 & 0.012856 \\
B & 7.5 & 0.005534 & 0.005663 & 0.006142 \\
B & 15 & 0.001950 & 0.002002 & 0.002620 \\
\bottomrule\end{tabular}
\caption{2024 results at thresholds absent from the fitting query grid. All three algorithms use 30\% of training anchors and the same eligible test population within each mode.}\label{tab:interior}
\end{table}
\FloatBarrier
GBT also exposes a structural issue with an unconstrained threshold predictor: a larger rainfall threshold can produce a larger exceedance probability. Across the ten evaluated thresholds, the proportion of adjacent query pairs with increasing GBT probability is 0.0491\% at 30\% training and 0.0681\% at 100\% for A, compared with 0.5238\% and 0.6164\% for B. RF preserves the required ordering by construction.

Appendix~\ref{app:prediction} retains threshold-specific baseline skill, monthly skill and feature-importance diagnostics for the 15\% models.

\FloatBarrier
\subsection{Worked probability query and model choice}
Table~\ref{tab:workedquery} illustrates the required interface using archived comparison scores for station S06 at 07:10 SGT on 8 November 2024. A retrieves the twelve published history endpoints from 06:05 to 07:00 and predicts accumulation over 07:10--07:40. Changing only $q$ gives non-increasing probabilities. The realised half-hour was dry despite A assigning 53.63\% probability above 1~mm; this illustrates forecast uncertainty and is not a representative performance sample or a newly executed live query.
\begin{table}[htbp]\centering\small
\begin{tabular}{rrr}\toprule
Threshold $q$ (mm) & A RF probability (\%) & B RF probability (\%)\\\midrule
0 & 82.50 & 6.67 \\
1 & 53.63 & 2.66 \\
10 & 13.01 & 0.43 \\
20 & 4.04 & 0.12 \\
\bottomrule\end{tabular}
\caption{Archived 15\% RF scores for one common station-time example. Both columns predict $P(Y>q)$; the observed $Y$ was 0~mm.}\label{tab:workedquery}
\end{table}

The \textbf{15\% depth-8 A RF is the fully evaluated reference for queries with complete published history}. Increasing its fitting fraction gives no clear mean improvement. Appendix~\ref{sec:rfsensitivity} finds lower 2022 Brier for depth 10 (0.01664745 versus 0.01669750 at 15\%), but that validation-only candidate needs calibration and subsequent evaluation before replacing the reference. For calendar-only risk, historical frequency remains the preferred baseline at the evaluated thresholds. Fine-resolution heavy-rain calibration and chronological confirmation of validation-selected candidates are the next priorities.
